\subsection{The binary-leaf reduction through center order six}

\begin{lemma}[Proved: a binary act represented by a two-element right ideal]
\label{CP2:A-binary-act-two-element-right-ideal}
Let $M$ be a finite monoid admitting an idempotent $e$ with
$|eM|\leq2$. Every two-element right $M$-set without a global
fixed point is isomorphic to the right act $eM$. In particular,
its action category is contractible.
\end{lemma}
\begin{proof}
Such an act $S$ is transitive: if a point $s$ had a proper
nonempty orbit $sM\subset S$, then $sM=\{s\}$ would make it
globally fixed. Choose $s\in Se$. Then $se=s$, and
\[
 eM\longrightarrow S,\qquad em\longmapsto sm
\]
is well-defined and right-$M$ equivariant: if $em=en$, then
$sm=sem=sen=sn$. Transitivity makes it onto. The cardinalities
force $|eM|=2$ and make the map a bijection. The right act $eM$
is a retract of the regular right act $M$, through left
multiplication by $e$. This induces an actual categorical
retraction of action categories. The regular action category
is contractible, so its retract is contractible as well.
\end{proof}

\begin{theorem}[Proved: binary trivial leaves through center order six]
\label{CP2:A-six-center-binary-trivial-leaf-nineteen}
A source star of at most nineteen arrows, with center order
at most six, one trivial leaf on a two-element stalk, and
other leaf order at most eleven, cannot realize $\CP^2$.
The opposite statement also holds.
\end{theorem}
\begin{proof}
Write $M$ for the center, $S$ for the binary stalk at the
trivial leaf, and $N,T$ for the other leaf and stalk.
A fixed point of $S$ yields the same proper cut retraction
as in Theorem~\ref{CP2:A-five-center-binary-trivial-leaf-nineteen}:
collapse $S$ to that point and fix $T$. Its image has at most
eighteen arrows and would retain actual projective-plane
type, contrary to Theorem~\ref{CP2:A-global-lower-nineteen}.
Assume henceforth that $S$ has no fixed point.

Let the minimal ideal have $r$ rows, $c$ columns, and local
group $G$, of order $g$. If $r=1$, the trivial left rational
module is projective. The actual exact-orbit quotient at
the center gives a rationally equivalent star with trivial
center. Its two ordinary act-cone wedge summands have
bounded rational homology and leaf orders at most eleven.
The bounded-eleven act theorem kills every second product,
as in the one-row case of the preceding theorem.

If $c=1$ and $g=1$, a minimal idempotent satisfies
$eM=eMe=\{e\}$. Lemma~\ref{CP2:A-small-minimal-ideal}
gives $BM\simeq *$. The full star is consequently an actual
wedge of its edge spaces. Projective-plane type of that
wedge forces an actual projective-plane edge summand, by
the split integral homology and primitive-square argument.
Each edge has at most seventeen arrows, a contradiction.

In every remaining case $r\geq2$ and either $c\geq2$ or
$g\geq2$. The minimal ideal is proper unless $M$ is a
group, which was already covered by $r=1$. Thus
\[
 rcg+1\leq |M|\leq6.
\]
If $c\geq2$, this forces $r=c=2$ and $g=1$.
If $c=1$, it forces $r=g=2$. In either case a minimal
right ideal $eM$ has $cg=2$ elements.
The preceding lemma makes $B(S,M)$ contractible.
The binary edge is the homotopy pushout
$BM\leftarrow B(S,M)\to *$, and its inclusion of $BM$
is a homotopy equivalence. Gluing the other edge along
$BM$ shows that deleting this binary edge and its trivial
leaf preserves the full star homotopy type. This removes
three arrows, leaving at most sixteen, again impossible.
\end{proof}
